%%%%%%%%%%%%%%%%%%%%%%%%%%%%%%%%%%%%%%%%%%%%%%%%%%%%%%%%%%%%%%%%%%%%%%%%%%%%%
%  enslyonstage-gabarit.tex — ENS de Lyon Biosciences internship report (blank template)
%
%  → Fill in the fields below, then write your sections.
%  → For a full worked example, see  enslyonstage-exemple.tex
%
%  On Overleaf: this file is detected automatically as the root document.
%  Locally: pdflatex enslyonstage-gabarit  →  biber enslyonstage-gabarit  →  pdflatex
%
%  Project structure:
%    enslyonstage-gabarit.tex ← this file  (fill in and compile)
%    enslyonstage-exemple.tex ← full worked example (for reference)
%    enslyonstage.cls         ← class file (same directory or on TEXMFHOME)
%    references.bib       ← BibTeX bibliography
%    figures/             ← images (jpg, pdf, png, eps …)
%    logos/               ← hautdepage.pdf  basdepage.pdf  Master.jpg
%                            logo_ens.pdf  (+ logo_labo.* optional)
%
%  Final PDF file name (ENS requirement):
%    LastName-FirstName-ENS-M1.pdf   or   LastName-FirstName-ENS-M2.pdf
%
% ═══════════════════════════════════════════════════════════════════════════
%  ENS DE LYON — OFFICIAL GUIDELINES  (keep these comments during writing)
% ═══════════════════════════════════════════════════════════════════════════
%  FORMAT
%    ✓ Written in English
%    ✓ Two-column layout (this class)        ✓ Pages numbered (automatic)
%    ✓ PDF, preferably < 4–5 MB
%    ✓ File name: "LastName-FirstName - ENS-M1.pdf"  or  "...-M2.pdf"
%    ✗ No supplemental data / annexes
%      (exceptions: contact stages.biosciences@ens-lyon.fr)
%
%  LENGTH & DISPLAY ITEMS
%    M1 : ≤ 5 000 words (± 10 %)   ≤ 6 figures + tables
%    M2 : ≤ 4 000 words (± 10 %)   ≤ 5 figures + tables
%    → word count excludes references and figure/table legends
%    → on Overleaf: Menu → Word Count  (uses the %TC: directives in the class)
%
%  TITLE     ≤ 125 characters including spaces
%  ABSTRACT  150–200 words · English · NO references
%  REFS      ≤ 35 publications · PNAS (new) numbered style [1],[2]…
%            Unpublished data / personal comm.: in-text only, NOT in ref list
%
%  STYLE
%    • Spell out 1–9 unless with units  ("three replicates" / "3 µm", "5 mM")
%    • Units: metric, negative exponents  (µmol m⁻² s⁻¹  NOT  µmol/m²/s)
%    • Gene names: italics  —  Protein names: roman  (follow field conventions)
%    • Abbreviations: define at first use, keep to a minimum
%    • Stats: specify test, n, and exact p-value for every numerical result
% ═══════════════════════════════════════════════════════════════════════════

% ── THEME ──────────────────────────────────────────────────────────────────
% Available (22): orange (default) · violet · blue · red · green · teal · gray
%                 navy · azure · indigo · midnight · slate · purple · pink
%                 crimson · wine · amber · gold · brown · lime · forest · cyan
\documentclass[orange]{enslyonstage}

% ── BIBLIOGRAPHY ───────────────────────────────────────────────────────────
% biblatex is compiled with biber (automatic on Overleaf)
\addbibresource{references.bib}

% ── HYPHENATION (optional) ─────────────────────────────────────────────────
% Add words that LaTeX hyphenates badly; use - to mark allowed break points.
% Example: \hyphenation{hy-phen-a-tion quad-ru-pe-dal}
%\hyphenation{}

% ── CUSTOM UNIT SHORTCUTS (optional) ───────────────────────────────────────
% Define frequently-used units as shortcuts to save typing.
% Add your own following the same pattern.
\newcommand{\um}{\si{\micro\metre}}       % \um   → µm
\newcommand{\uM}{\si{\micro\Molar}}       % \uM   → µM
\newcommand{\degC}{\si{\degreeCelsius}}   % \degC → °C

%%%%%%%%%%%%%%%%%%%%%%%%%%%%%%%%%%%%%%%%%%%%%%%%%%%%%%%%%%%%%%%%%%%%%%%%%%%%%
%  METADATA  —  fill in every field below
%%%%%%%%%%%%%%%%%%%%%%%%%%%%%%%%%%%%%%%%%%%%%%%%%%%%%%%%%%%%%%%%%%%%%%%%%%%%%

% ── TITLE ──────────────────────────────────────────────────────────────────
% ENS requirement: maximum 125 characters including spaces.
% To force a line break in the title without breaking PDF metadata, use:
%   \titre{First part:\texorpdfstring{\\}{ }Second part}
\titre{}

% ── AUTHOR ─────────────────────────────────────────────────────────────────
% Use \textsc{} for small caps on last name: Firstname \textsc{Lastname}
\auteur{\textsc{}}

% ── DEGREE LEVEL ───────────────────────────────────────────────────────────
% Enter "M1" or "M2" (without "Biosciences" — added automatically)
\master{M}

% ── BIOLOGICAL DOMAIN ──────────────────────────────────────────────────────
% Shown on cover page and in footer.
% Examples: Molecular and Cellular Biology, Neurosciences, Ecology, etc.
\domaine{}

% ── INTERNSHIP DATES ───────────────────────────────────────────────────────
% Format: DD Month YYYY (e.g., 02 February 2025)
\datedebut{}
\datefin{}

% ── LABORATORY ─────────────────────────────────────────────────────────────
% Full laboratory name, institute, and location.
% Example: Lab Name, Institute Name, CNRS UMR XXXX, City
\laboratoire{%

    }

% ── SUPERVISORS ────────────────────────────────────────────────────────────
% Lab supervisor (internship mentor) — name and email
\tuteurstage{\textsc{}}{}
% ENS academic supervisor — name and email
\tuteurENS{\textsc{}}{}

% ── ABSTRACT ───────────────────────────────────────────────────────────────
% ENS requirements:
%   - 150–200 words (single paragraph)
%   - English language
%   - NO literature references
%   - Include: (1) brief background and question, (2) main results without
%     extensive detail, (3) summary of significance
%   - Keep accessible to non-specialists
%
% Word count: the abstract is NOT counted toward the ENS word limit.
% (This is handled automatically by the %TC: directives in enslyonstage.cls.)
\abstracts{%

}

% ── KEYWORDS ───────────────────────────────────────────────────────────────
% ENS guideline: 5 to 8 keywords, comma-separated
\keywords{}

% ── WORD COUNT ─────────────────────────────────────────────────────────────
% ENS limits: M1 ≤ 5000 words (±10%)  /  M2 ≤ 4000 words (±10%)
% Count excludes abstract, references, and figure/table legends.
%
% On Overleaf: Menu → Word Count
%   The class automatically excludes figure captions and table environments
%   via %TC: directives — no extra setup needed.
%
% Locally: texcount -sum -1 enslyonstage-gabarit.tex
%
% Update the number below before each submission:
\nombredemots{}

% ── PREVIOUS INTERNSHIPS ───────────────────────────────────────────────────
% Table of previous internships shown on cover page.
% Use \stagerow{Year}{Lab, City, Country}{Supervisor}{Topic}
% Leave cells empty for levels not yet completed.
\stagesanterieurs{
	\stagerow{L3}{}{\textsc{}}{}
	\stagerow{M1}{}{\textsc{}}{}
	\stagerow{M2}{}{\textsc{}}{}
}

%%%%%%%%%%%%%%%%%%%%%%%%%%%%%%%%%%%%%%%%%%%%%%%%%%%%%%%%%%%%%%%%%%%%%%%%%%%%%
%  LOGOS  —  place files in the logos/ directory
%  The cover page uses hautdepage.pdf and basdepage.pdf automatically.
%%%%%%%%%%%%%%%%%%%%%%%%%%%%%%%%%%%%%%%%%%%%%%%%%%%%%%%%%%%%%%%%%%%%%%%%%%%%%

% ── ENS LOGO — displayed in footer center (automatic)
\logoENS[1.5cm]{logos/logo_ens.pdf}

% ── MASTER LOGO — kept for compatibility
\logoformation[2cm]{logos/Master.jpg}

% ── LABORATORY LOGO — shown on cover page (optional, remove if not needed)
% \logoLabo[2.5cm]{logos/logo_labo.pdf}

%%%%%%%%%%%%%%%%%%%%%%%%%%%%%%%%%%%%%%%%%%%%%%%%%%%%%%%%%%%%%%%%%%%%%%%%%%%%%
%  HEADER / FOOTER CUSTOMIZATION (optional)
%  Defaults work well for most reports — uncomment only if needed.
%%%%%%%%%%%%%%%%%%%%%%%%%%%%%%%%%%%%%%%%%%%%%%%%%%%%%%%%%%%%%%%%%%%%%%%%%%%%%
% Default layout:
%   top-left  : current section name (\leftmark)
%   top-right : title wrapped at 0.6\linewidth
%   bot-left  : \domaine in small caps
%   bot-centre: \logoENS at 8mm (below footer rule)
%   bot-right : page/total

%\coinhautgauche{custom text}
%\coinhautdroit{custom text}
%\coinbasgauche{custom text}
%\coinbascentre{\includegraphics[height=8mm]{logos/custom.pdf}}
%\coinbasdroit{custom text}

%%%%%%%%%%%%%%%%%%%%%%%%%%%%%%%%%%%%%%%%%%%%%%%%%%%%%%%%%%%%%%%%%%%%%%%%%%%%%
\begin{document}
%%%%%%%%%%%%%%%%%%%%%%%%%%%%%%%%%%%%%%%%%%%%%%%%%%%%%%%%%%%%%%%%%%%%%%%%%%%%%

% ── COVER PAGE — generates ENS cover, then switches to two-column mode
\fairepagedegarde

% ── ABSTRACT BOX — full-width colored box; must follow \fairepagedegarde
\faireabstract

% ── TABLE OF CONTENTS — optional; remove if your tutor does not require it
%\tableofcontents

%%%%%%%%%%%%%%%%%%%%%%%%%%%%%%%%%%%%%%%%%%%%%%%%%%%%%%%%%%%
\section{Introduction}
\label{sec:intro}
%%%%%%%%%%%%%%%%%%%%%%%%%%%%%%%%%%%%%%%%%%%%%%%%%%%%%%%%%%%
% ENS requirements:
%   - Background with MORE detail than a journal paper
%   - Understandable to non-specialists
%   - End with a clear statement of the question(s) addressed
%   - Define non-standard abbreviations at first use
%
% TIP — drop cap on the opening letter:
%   \lettrine[lines=3]{T}{he} context of this study...
%   (lettrine is already loaded by enslyonstage.cls)

\paragraph{}


% \newpage   ← uncomment to force a column/page break if needed

%%%%%%%%%%%%%%%%%%%%%%%%%%%%%%%%%%%%%%%%%%%%%%%%%%%%%%%%%%%
\section{Materials and Methods}
\label{sec:methods}
%%%%%%%%%%%%%%%%%%%%%%%%%%%%%%%%%%%%%%%%%%%%%%%%%%%%%%%%%%%
% ENS requirements:
%   - Sufficient detail to allow repetition
%   - Standard procedures may be referenced
%   - Statistical tests: name, n, and exact p-value for every result
%
% Use \paragraph{Heading.} for each experimental approach.

\paragraph{}


%%%%%%%%%%%%%%%%%%%%%%%%%%%%%%%%%%%%%%%%%%%%%%%%%%%%%%%%%%%
\section{Results}
\label{sec:results}
%%%%%%%%%%%%%%%%%%%%%%%%%%%%%%%%%%%%%%%%%%%%%%%%%%%%%%%%%%%
% ENS requirements:
%   - Findings presented clearly and succinctly
%   - Concise summary of data in ALL display items
%   - Display item limits: M1 ≤ 6 figures+tables  /  M2 ≤ 5 figures+tables
%   - Legends must be self-sufficient
%
% Figure insertion:
%   \insererfigure[options]{filepath}{caption text}{fig:label}
%   Options: col (single-column)  p (full-page)  b (bottom)
%            h=<dim> (height)     w=<dim> (width)
%   Default: double-column, h=0.33\textheight
%   See enslyonstage-exemple.tex for worked examples of every option.
%
% Multi-reference coloured cross-reference (cleveref, already loaded):
%   \crefcolor{fig:a,fig:b}   → coloured "Figs. 1 and 2"
%   Do NOT add \usepackage{cleveref} — it is already loaded by the class.

\subsection{}


%%%%%%%%%%%%%%%%%%%%%%%%%%%%%%%%%%%%%%%%%%%%%%%%%%%%%%%%%%%
\section{Discussion}
\label{sec:discussion}
%%%%%%%%%%%%%%%%%%%%%%%%%%%%%%%%%%%%%%%%%%%%%%%%%%%%%%%%%%%
% ENS requirements:
%   - Significance of results in broad biological context
%   - Critical evaluation — do NOT oversell
%   - Shortcomings and how they could be addressed
%   - 2–3 concrete follow-up experiments
%
% This is NOT a journal manuscript — show your capacity to critically
% discuss your work, including negative results and limitations.

\paragraph{}


%%%%%%%%%%%%%%%%%%%%%%%%%%%%%%%%%%%%%%%%%%%%%%%%%%%%%%%%%%%
\section*{Author Contributions}
%%%%%%%%%%%%%%%%%%%%%%%%%%%%%%%%%%%%%%%%%%%%%%%%%%%%%%%%%%%
% ENS requirement (mandatory):
%   - Your personal contribution (which experiments/analyses you did)
%   - Role of EVERY person who participated in your work
%
% Format: "F.L. designed and performed experiments X and Y;
%          J.S. conceived the project and supervised the internship."

\paragraph{}


%%%%%%%%%%%%%%%%%%%%%%%%%%%%%%%%%%%%%%%%%%%%%%%%%%%%%%%%%%%
\section*{Acknowledgements}
%%%%%%%%%%%%%%%%%%%%%%%%%%%%%%%%%%%%%%%%%%%%%%%%%%%%%%%%%%%
% Optional — thank people not listed in Author Contributions
% (technical staff, facility managers, funding sources, etc.)

This report was typeset using a custom \LaTeX{} class developed for ENS de Lyon Biosciences internship reports.

%--- Bibliography — ≤ 35 references (ENS guideline) --------------------------
% Balance columns on the final page (cleaner last page)
\balance
\biblio

\end{document}
